\section{Limitations}
\label{sec:limitations}

\paragraph{Identification at scale does not work.} At the measured code distribution there is no
parameter setting that is simultaneously secure and usable at $N = 5{,}000$
(\S\ref{sec:params}). The design is coherent as a verification $(1{:}1)$ system and as a
small-$N$ identification system, and we do not claim more.

\paragraph{The embedding is the ceiling.} Every headline number here would move with a
higher-dimensional or better-trained extractor, and none would move with a better coder or a
different threshold.

\paragraph{Templates at rest.} The stored code is a bearer credential and fails all three
ISO/IEC 24745 requirements (\S\ref{sec:atrest}). Our mitigations recover unlinkability and
renewability but not irreversibility. This posture is \emph{weaker} than the encrypted-template
baseline's, and we treat full at-rest confidentiality against a compromised server as out of
scope rather than solved.

\paragraph{Semi-honest only.} A malicious client can deviate arbitrarily, and since it supplies
$\xt$ with no binding to a live finger, sensor attestation is load-bearing rather than a detail.

\paragraph{Evaluation scope.} All results are on a single corpus whose probes are synthetic
degradations --- obliteration, rotation, z-cut --- of the \emph{same} capture. That measures
robustness to noise, not generalisation across capture conditions, sensors or sessions, which
remains untested here and is the first thing further work should establish.

\paragraph{Communication.} Per-query communication exceeds the homomorphic baseline's beyond a
few hundred records and grows linearly; the design trades communication for storage and key
material.
